\section{Discussion}

Our results suggest that the limitations of current agents cannot be addressed through a single optimization strategy.
Different failure patterns require corresponding changes to model training, inference-time strategies, long-horizon system design, or the evaluation objective itself.

\textbf{What Training Can Improve}

Our process analysis suggests that training should be tailored to the specific weaknesses of each model and task category.
Execution is already strong and tightly clustered across models, so generic code-execution training alone is unlikely to be the primary field-wide opportunity.
Solution Framing and Feedback Control vary more widely, indicating greater scope for model-specific improvements in direction selection and feedback use.
Training priorities should likewise vary across task categories according to whether the main bottleneck lies in framing, implementation, or feedback control.
These process metrics can guide the construction of targeted training data, process rewards, and curricula beyond what final reward alone provides.
The experience experiments offer an additional training signal: paired cases of positive and negative transfer could help models learn when prior experience is applicable and when it should be reconsidered.

\textbf{What Inference-Time Search Can Recover}

Beyond training, inference-time strategies offer a direct way to improve how reliably models realize their existing capabilities.
The contrast between average and best-run performance shows that several models can reach competitive solutions but do not reproduce them consistently.
This creates an opportunity to generate more diverse rollouts and use verifier feedback to identify promising trajectories.
Instead of assigning every trajectory a fixed budget, compute could be redirected by branching from promising checkpoints and terminating trajectories that repeatedly fail or stagnate.
Process diagnostics could further guide this allocation by encouraging broader exploration when Solution Framing is weak and deeper implementation or recovery when Execution or Feedback Control is limiting.
Trajectory selection could combine verifier outcomes with execution validity and progress retention to avoid relying on final reward alone.

\textbf{What Memory and Harness Design Can Stabilize}

Some failures arise from retaining and applying information over a long research process rather than from generating an effective action in isolation.
Accumulated experience usually helps agents preserve useful discoveries and avoid known failures, but it can also carry forward misleading conclusions or anchor exploration to a local optimum.
An effective memory system therefore requires more than simply storing additional context. It must support the selective retrieval, validation, revision, and removal of experience according to the current task.

Harnesses address a complementary aspect of long-horizon control.
Native harnesses improve run-to-run stability but do not substantially raise the observed performance ceiling.
This stability gain likely comes from mechanisms that reduce avoidable failures during extended experimentation, including error recovery, task management, and best-state protection.
Automated harness optimization offers further headroom, including task-specific harnesses that target distinct research bottlenecks and model-adaptive harnesses that account for differences in tool use, planning, and recovery behavior.

\textbf{What Requires New Objectives and Benchmarks}

Some limitations cannot be resolved through training, inference-time strategies, memory, or harness design when the reward captures task performance but not methodological quality.
Current agents execute and optimize effectively, yet their strongest solutions primarily compose established techniques, while validated novel approaches remain rare.
Moreover, evaluator-specific shortcuts are substantially more common than novel approaches when agents depart from standard solutions.
More aggressive optimization of the same reward may therefore reinforce shortcut-seeking rather than improve research quality.
Progress toward more open-ended research and scientific discovery will require tasks and feedback that reward not only task performance, but also novelty, validity, and generality.






